% mn2esample.tex
%
% v2.1 released 22nd May 2002 (G. Hutton)
%
% The mnsample.tex file has been amended to highlight
% the proper use of LaTeX2e code with the class file
% and using natbib cross-referencing. These changes
% do not reflect the original paper by A. V. Raveendran.
%
% Previous versions of this sample document were
% compatible with the LaTeX 2.09 style file mn.sty
% v1.2 released 5th September 1994 (M. Reed)
% v1.1 released 18th July 1994
% v1.0 released 28th January 1994

\documentclass[useAMS,usenatbib]{mn2e}
\usepackage{amsmath,graphicx,xspace,hyperref,url,acronym,aas_macros}


% If your system does not have the AMS fonts version 2.0 installed, then
% remove the useAMS option.
%
% useAMS allows you to obtain upright Greek characters.
% e.g. \umu, \upi etc.  See the section on "Upright Greek characters" in
% this guide for further information.
%
% If you are using AMS 2.0 fonts, bold math letters/symbols are available
% at a larger range of sizes for NFSS release 1 and 2 (using \boldmath or
% preferably \bmath).
%
% The usenatbib command allows the use of Patrick Daly's natbib.sty for
% cross-referencing.
%
% If you wish to typeset the paper in Times font (if you do not have the
% PostScript Type 1 Computer Modern fonts you will need to do this to get
% smoother fonts in a PDF file) then uncomment the next line
% \usepackage{Times}

%%%%% AUTHORS - PLACE YOUR OWN MACROS HERE %%%%%
% Macros
\newcommand{\mgii}{Mg~{\sc ii}\xspace}

\newcommand{\ntcomment}[1]{\texttt{\\ \noindent[NT: #1]}}


% Acronyms
\acrodef{los}[LOS]{line-of-sight}
\acrodef{lss}[LSS]{large scale structure}
\acrodef{sdss}[SDSS]{Sloan Digital Sky Survey}
\acrodef{redmapper}[redMaPPer]{red-sequence Matched-filter Probabilistic Percolation}
\acrodef{qso}[QSO]{quasi-stellar object}
\acrodef{snr}[SNR]{signal-to-noise ratio}
%%%%%%%%%%%%%%%%%%%%%%%%%%%%%%%%%%%%%%%%%%%%%%%%

\title[The incidence of \mgii\ as a function of cluster masses]{Cold
  gas within and around galaxy clusters: The incidence of \mgii\ as a
  function of cluster masses \ntcomment{better title TBD}}
\author[Crighton, Tejos \& L\'opez]{
\parbox[t]{\textwidth}{
\vspace{-1.0cm} 
Neil H. Crighton$^{1}$\thanks{E-mail: \href{mailto:neilcrighton@gmail.com}{neilcrighton@gmail.com}},
Nicolas Tejos$^{2}$ and
Sebasti\'an L\'opez$^{3}$
}
\vspace*{6pt} \\ 
$^{1}$ Centre for Astrophysics and Supercomputing, Swinburne University of Technology, Hawthorn, Victoria 3122, Australia\\
$^{2}$ Department of Astronomy and Astrophysics, UCO/Lick Observatory, University of California, 1156 High Street, Santa Cruz, CA 95064, USA\\
$^{3}$ Departamento de Astronom\'ia, Universidad de Chile, Casilla 36-D, Santiago, Chile\\
\vspace*{-0.5cm}}

\begin{document}

\date{Draft v1.0}

\pagerange{\pageref{firstpage}--\pageref{lastpage}} \pubyear{2014}

\maketitle

\label{firstpage}

\begin{abstract}
We investigate the incidence of \mgii\ absorption line systems in
galaxy clusters as a function of cluster mass. Our results are....
\end{abstract}

\begin{keywords}
--intergalactic medium --quasars: absorption lines --galaxies: clusters:
intracluster medium --large scale structure of the Universe --galaxies:
evolution
\end{keywords}

\section{Introduction}
\acresetall
\ntcomment{brief introduction}

\section{Data}

\subsection{Galaxy clusters}
We use the \ac{redmapper} catalog \citep{Rykoff2014} ...
\ntcomment{\\
- Description of the \ac{redmapper} cluster catalog\\
- Explain why \ac{redmapper}: advantages\\
- Clusters centers as probabilities\\
- Cluster members as probabilities\\
- Cluster Richness corrected by incompleteness\\
- Photometric uncertainties are small\\
- Some spectroscopic redshift exists\\
- Adopted richness--mass relationship\\
- Explain completeness and purity\\}

\subsection{\mgii absorption line systems}
We use the \mgii\ catalog published by  ...

\ntcomment{\\
- Description of the \mgii\ catalog\\
- Explain why this catalog in particular: advantages\\
- Explain the main observables relevant for this study: z, EW, Vdisp? \\
- Explain why completeness level issues, and why these do not affect our analysis.}

\section{Data analysis and results}
\subsection{Cross-match between clusters and \mgii\ systems}
\ntcomment{\\
- Show plots of redshift overlap\\
- Explain subsample of the cluster catalog used in this analysis\\
- Explain subsample of the \mgii\ catalog used in this analysis\\
- Justify the maximum scale adopted (~40 Mpc? use as motivation the mean distance between clusters?)
- Define redshift path per cluster\\
- Define hits\\
}
\begin{figure}
 \begin{minipage}{0.5\textwidth}
    \centering
    \includegraphics[width=5cm]{figs/figure_sample.pdf}
  \end{minipage}

\caption{Clusters and \mgii\ redshift distribution...}
\label{fig:fig1}

\end{figure}


\subsection{The incidence \mgii\ in clusters}
\ntcomment{\\
- Define incidence in clusters, $dN/dz$\\
- Define incidence as a function of co-moving separation\\
- Define incidence as a function of virial radii; emphasize why is important to scale by virial mass: massive clusters are larger. Dark matter haloes are expected to show self-similarity, others)\\
- Plot these two quantities for our full samples in both co-moving and r200 separations.\\
- Mention whether these two seem consistent with each other but leave interpretation to the Discussion section. If no qualitative differences, stick to r200.}

\begin{figure}
 \begin{minipage}{0.5\textwidth}
    \centering
    \includegraphics[width=5cm]{figs/figure_sample.pdf}
  \end{minipage}

\caption{$dN/dz$ for our full sample; in both, co-moving and r200 units}
\label{fig:fig2}

\end{figure}


\subsection{The incidence \mgii\ in clusters as a function of cluster mass}
\ntcomment{\\
- Define cluster mass bins; explain why we chose these particular bins.\\
- Main result of the paper: incidence as a function of cluster masses\\
- Plot incidence as a function of cluster masses, in r200 (or both in case is worth it)}



\subsection{Physical model}
\ntcomment{\\ 
- Explain our model to fit the results\\ 
- Mention whether our model consistent with the known masses of the
clusters in each mass bin?  (either result will be very
interesting) but leave full interpretation to Discussion section\\ }

\begin{figure}
 \begin{minipage}{0.5\textwidth}
    \centering
    \includegraphics[width=5cm]{figs/figure_sample.pdf}
  \end{minipage}

\caption{$dN/dz$ for our cluster-mass subsamples; in both, co-moving and r200 units; this figure should contain our model}
\label{fig:fig3}

\end{figure}


\section{Discussion}
\ntcomment{\\
- General implications of results, focusing on the main result\\
- Although we focused on the larger scales, we can speculate now on our reported results on scales < 1 Mpc (subject to large statistical uncertainties).\\ 
- Why is important to have cold gas in clusters\\
- Are the observed trends consistent with expectations?\\
}

\subsection{Comparison to previous work}
\ntcomment{\\
- Lopez+2008\\
- Zhu+14\\
- Gauthier+14\\
- etc. \\
- emphasize why our experiment is different (we focus on the most massive haloes in the universe and we know the masses of the clusters very well)\\
- focus on agreements and disagreements\\
}

\subsection{Future prospects [??]}
\ntcomment{\\
- What can be done to extend this work?\\
- Is it worth to aim for less massive groups? Is there evidence for a different behaviour in these lower mass haloes?\\
- Mention work on the smaller scales that we will do ?\\
- Address different galaxy populations (star-forming, non-star-forming)\\
- Split by \mgii\ EW for a fixed cluster mass\\
- Investigate properties of clusters showing and not showing \mgii.\\
- Put limits on the covering fraction of \mgii inside galaxy clusters [this seems important, should we address this in this paper rather than later? We could to it for the strongest absorbers without worrying much about incompleteness].\\
}



\section{Summary}
\ntcomment{Brief summary}

\section*{Acknowledgments}

We thank contributors to SciPy\footnote{\url{http://www.scipy.org}},
Matplotlib\footnote{\url{http://www.matplotlib.sourceforge.net}},
Astropy\footnote{\url{http://www.astropy.org}\citep{AstropyCollaboration2013}},
and the Python programming
language\footnote{\url{http://www.python.org}}; the free and
open-source community; and the NASA Astrophysics Data
System\footnote{\url{http://adswww.harvard.edu}} for software and
services.

Add \ac{sdss} thanks.

N.H.C. acknowledges...; N.T. acknowledges...; S.L. acknowledges...

\bibliographystyle{mn2e_trunc8.bst}
\bibliography{/home/ntejos/lit/bib/IGM}

\appendix

\section{Check for systematic effects}
\ntcomment{\\ 
-Plot QSO properties as a function of cluster masses, and transverse
separation in the same way the analysis was presented.\\
-Plot properties of clusters (mass distribution, redshift distribution)
with \mgii\ versus those without ?}
\begin{figure}
 \begin{minipage}{0.5\textwidth}
    \centering
    \includegraphics[width=5cm]{figs/figure_sample.pdf}
  \end{minipage}

\caption{Properties of \acp{qso} as a function of transverse distance.}
\label{fig:fig_appendix}

\end{figure}



\bsp
\label{lastpage}
\end{document}
